\section{Dark crossings in exact symmetry quotients}
\label{sec:rewrite-crossings}

If $\Pi_a$ and $\Pi_b$ project onto inequivalent irreps of a symmetry
preserved by the entire path, then
\begin{equation}
 \Pi_a\,\partial_sH(s)\,\Pi_b=0.
 \label{eq:dark-coupling-rewrite}
\end{equation}
An equality of the two branch energies is therefore a global crossing but
not an avoided crossing seen by a state confined to either block.

This phenomenon is stable within the symmetry-preserving class; it is not
restricted to exactly solvable diagonal cancellation points.

\begin{theorem}[Stability of a transverse dark crossing]
\label{thm:dark-crossing-stability}
Let $H_\varepsilon(s)=H(s)+\varepsilon R(s)$ be a continuously
differentiable finite-dimensional path, where $\varepsilon\in\mathbb R$ is
the perturbation strength and $R(s)$ is Hermitian.  Let $\Pi_a,\Pi_b$
project onto
inequivalent symmetry sectors.  Suppose every $H_\varepsilon(s)$ commutes
with both projectors.  Assume that near $(s_\star,0)$ the two restricted
blocks have simple isolated eigenvalue branches $E_a(s,\varepsilon)$ and
$E_b(s,\varepsilon)$ such that
\begin{equation}
 E_a(s_\star,0)=E_b(s_\star,0),
 \qquad
 \partial_s(E_a-E_b)(s_\star,0)\ne0.
 \label{eq:transverse-dark-hypothesis}
\end{equation}
Then, for all sufficiently small $|\varepsilon|$, there is a unique nearby
$s_\star(\varepsilon)=s_\star+O(|\varepsilon|)$ at which the two branches
cross exactly, and
\begin{equation}
 \Pi_a\partial_sH_\varepsilon(s)\Pi_b=0.
\end{equation}
If the $a$ branch is separated from the rest of its block by at least
$\Delta_\star>0$ throughout a neighbourhood of $s_\star$ at
$\varepsilon=0$, then its isolation gap remains at least
$\Delta_\star/2$ whenever
$|\varepsilon|\sup_{s\in[0,1]}\|R(s)\|\leq\Delta_\star/4$ and the perturbed crossing
stays in that neighbourhood.
\end{theorem}

Physically, the preserved symmetry forbids the two branches from
hybridizing.  A symmetry-preserving perturbation may move the crossing, but
cannot open an avoided crossing between the sectors.  The complete
stability argument is in Appendix~\ref{app:dark-stability}.

\subsection{A finite joint-symmetry audit}

For a frozen $2\times4$ incidence instance with $k=3$, $\lambda=5$, and
$\mu=10$,
$G_B\simeq D_4$, where $D_4$ is the order-eight dihedral group, and the
$S_3\times D_4$ fixed quotient has dimension $22$,
compared with $512$ in the valid ordered space.  Cyclic closure from
$\ket u$ also has dimension $22$, as determined independently by the
generator algebra.  The uncatalysed path has a fixed-space rank closure at
$s=7/15$, showing that joint symmetry alone does not guarantee a rank-one
accessible gap.

As a finite illustration of a dark crossing away from the cancellation
point, use the single-register walk $W$ from
Eq.~\eqref{eq:walk-definition} and define
\begin{equation}
 C_2=-\sum_{1\leq r<t\leq k}W^{(r)}W^{(t)},
 \label{eq:d4-catalyst-operator}
\end{equation}
and use the endpoint-preserving path
\begin{equation}
 H_\chi(s)=H(s)+4\chi s(1-s)C_2,
 \label{eq:d4-catalyst-path}
\end{equation}
where $\chi\in\mathbb R$ is the dimensionless catalyst strength.
The operator $C_2$ commutes with both $S_k$ and $G_B$, and its envelope
vanishes at $s=0,1$.  At $\chi=2$, a crossing between different $D_4$
irreps occurs at $s\simeq0.806454$, while the $D_4$-trivial gap there is
approximately $0.258$.  This diagnostic is not used in the asymptotic
theorem.  Appendix~\ref{app:d4-catalyst} specifies its numerical protocol,
minimum-gap scan, and residuals.

\subsection{An exact multiplicity closure on the even-cycle family}

The asymptotic witness is analytic.  Express minimum vertex cover on the
even cycle $C_n$, $n=2k$, as set cover: bases are vertices, tasks are edges,
and each base covers its two incident edges.  The only minimum covers are
the two alternating subsets, which form one orbit of $D_n$.
Here $D_n$ denotes the dihedral group generated by the rotations and
reflections of $C_n$.

For this family $n_B=n$.  Consider the original ordered-register linear path
$H(s)$ in Eq.~\eqref{eq:cycle-linear-rewrite}.
The off-diagonal coefficient of each single-register complete-graph move is
\begin{equation}
 -\frac{1-s}{n}+\frac{s}{n-1},
\end{equation}
and vanishes exactly at
\begin{equation}
 s_c=\frac{n-1}{2n-1}.
 \label{eq:cycle-cancel-rewrite}
\end{equation}
At this point the Hamiltonian is diagonal up to a scalar.  All $2k!$
ordered representatives of the two alternating covers are global ground
states, whereas their normalized orbit is the unique ground vector in the
$S_k\times D_n$ fixed space.  Assuming $\mu\geq\lambda$, every other orbit
has diagonal penalty at least $\lambda$.  This lower bound is attained: the
distinct-base subset
\begin{equation}
 S_\star=\{1\}\cup\{2j:1\leq j\leq k-1\}
\end{equation}
has exactly one uncovered edge and no duplicate.  Its joint orbit state
therefore lies exactly $\lambda s_c$ above the ground energy.  Write
$g_0^{\rm global}$ and $g_0^{\rm sym}$ for the ranks of the global and
joint-fixed ground projectors, and write $\Delta_{\rm sym}$ for the
joint-fixed excitation gap $\Delta_{\rm exc}^{\mathcal H_{\rm sym}}$
defined in Section~\ref{sec:rewrite-gaps}.  Consequently,
\begin{align}
 g_0^{\rm global}(s_c)&=2k!,&
 g_0^{\rm sym}(s_c)&=1,\\
 \Delta_{\rm sym}(s_c)&=\lambda\frac{n-1}{2n-1}.
 \label{eq:cycle-dark-closure-rewrite}
\end{align}
Thus the global multiplicity gap closes exactly while the joint-fixed gap
is bounded below by a constant.  The multiplicity closure is dynamically
dark for the uniform initial state and the symmetry-preserving path.  No
transversality of the coincident branches is asserted here.

Equation~\eqref{eq:cycle-dark-closure-rewrite} controls the closure point,
not the minimum gap on every part of the linear path.  The next section
introduces a different, explicit parent path on the same family and proves
a uniform polynomial bound.  Keeping these statements separate avoids
turning a pointwise multiplicity-closure statement into an unsupported
adiabatic claim.
